% ============================================================
%  TracePay: IEEE Conference Paper
%  Compile with: pdflatex -> bibtex -> pdflatex -> pdflatex
% ============================================================
\documentclass[conference]{IEEEtran}

% ── Packages ────────────────────────────────────────────────
\usepackage[T1]{fontenc}
\usepackage[utf8]{inputenc}
\usepackage{cite}
\usepackage{amsmath,amssymb}
\usepackage{algorithmic}
\usepackage{algorithm}
\usepackage{array}
\usepackage{url}
\usepackage{hyperref}
\usepackage{graphicx}
\usepackage{booktabs}
\usepackage{listings}
\usepackage{xcolor}
\usepackage{balance}
\usepackage{microtype}

% ── Solidity listing style ──────────────────────────────────
\lstdefinelanguage{Solidity}{
  keywords={pragma,solidity,contract,interface,function,returns,
            external,public,internal,private,view,pure,payable,
            nonReentrant,override,virtual,event,emit,mapping,
            address,uint256,uint8,uint16,uint64,uint80,int256,
            bool,bytes32,string,struct,enum,constructor,require,
            revert,import,if,else,return,true,false,immutable,
            constant,memory,calldata,storage},
  sensitive=true,
  morecomment=[l]{//},
  morecomment=[s]{/*}{*/},
  morestring=[b]",
}
\lstset{
  language=Solidity,
  basicstyle=\ttfamily\scriptsize,
  keywordstyle=\color{blue!70!black}\bfseries,
  commentstyle=\color{green!50!black}\itshape,
  stringstyle=\color{red!60!black},
  breaklines=true,
  frame=single,
  numbers=left,
  numberstyle=\tiny\color{gray},
  xleftmargin=1.5em,
  framexleftmargin=1.5em,
  tabsize=4,
  captionpos=b,
}

% ── Metadata ────────────────────────────────────────────────
\hypersetup{
  pdftitle={TracePay: A Blockchain-Based Supply-Chain Trade Finance Platform with Chainlink Oracle Integration and Programmable Escrow},
  pdfauthor={Karen Elisha Chezhiyan},
  colorlinks=false,
}

% ============================================================
\begin{document}

\title{TracePay: A Blockchain-Based Supply-Chain Trade Finance Platform
       with Oracle-Priced Escrow, Milestone Verification,
       and On-Chain Invoice Financing}

\author{
  \IEEEauthorblockN{Karen Elisha Chezhiyan}
  \IEEEauthorblockA{\textit{School of Computer Science and Engineering}\\
    \textit{RV University}\\
    Bangalore, India\\
    karenelishachezhiyanbtech24@rvu.edu.in}
  \and
  \IEEEauthorblockN{Name}
  \IEEEauthorblockA{\textit{School of Computer Science and Engineering}\\
    \textit{RV University}\\
    Bangalore, India\\
    mail}
  \and
  \IEEEauthorblockN{Name}
  \IEEEauthorblockA{\textit{School of Computer Science and Engineering}\\
    \textit{RV University}\\
    Bangalore, India\\
    mail}
}

\maketitle

% ── Abstract ────────────────────────────────────────────────
\begin{abstract}
Global supply-chain trade finance remains plagued by information
asymmetry, counterparty risk, slow settlement cycles, and the high
cost of invoice-factoring intermediaries.
This paper presents \textit{TracePay}, a decentralized application
(DApp) that encodes a complete trade-finance lifecycle inside a single
Ethereum smart contract.
TracePay combines five tightly integrated mechanisms:
(1) keccak256 invoice hashing for duplicate-prevention and
off-chain privacy;
(2) Chainlink ETH/USD oracle-priced escrow with freshness validation;
(3) role-gated logistic milestone recording (Dispatched
$\rightarrow$ Delivered) with IPFS evidence CIDs;
(4) Delegation-of-Authority (DoA) multi-signature approval for
high-value orders exceeding a configurable USD threshold; and
(5) on-chain invoice factoring that grants suppliers immediate
working capital at a configurable advance ratio (default 90\%),
with interest settled atomically through a pull-payment waterfall.
The contract is implemented in Solidity 0.8.24 using
OpenZeppelin's \texttt{AccessControl} and \texttt{ReentrancyGuard}
primitives, and is validated by 16 Hardhat test cases (6 functional,
10 negative/revert) all passing with 100\% function coverage.
Gas benchmarks show deployment at 7.4\% of the Ethereum block limit,
and per-transaction costs ranging from 30\,628 to 181\,063 gas units.
\end{abstract}

\begin{IEEEkeywords}
blockchain, smart contract, trade finance, supply chain,
Ethereum, Chainlink oracle, escrow, invoice factoring,
decentralised finance (DeFi), role-based access control (RBAC),
delegation of authority, Solidity
\end{IEEEkeywords}

% ============================================================
\section{Introduction}
\label{sec:intro}

Trade finance is the financial backbone of international commerce,
bridging the gap between a buyer's need to defer payment and a
supplier's need for immediate liquidity \cite{icc2022}.
The traditional model relies on letters of credit, bank guarantees,
and invoice-factoring intermediaries, each introducing delays of
5--30\,days and fees of 2--8\% of the transaction value \cite{bis2019}.
Worse, 45\% of small and medium enterprises (SMEs) report rejected or
delayed trade-finance applications due to insufficient credit history
or documentation errors \cite{adb2021sme}.

Blockchain technology---and Ethereum smart contracts in particular---offer
a path toward trust-minimised, programmable settlement in which
payment is released automatically when verifiable conditions are met
\cite{nakamoto2008,buterin2014}.
Several platforms have explored this space
\cite{tracr2020,wetrade2022,contour2020}, yet most operate in
permissioned consortium settings and expose limited programmability.

We present \textbf{TracePay}, an open, permissionless DApp that:
\begin{itemize}
  \item Encodes the entire buyer--supplier--financier lifecycle in a
        single auditable smart contract ($\sim$220 SLOC of Solidity).
  \item Prices escrow in USD via live Chainlink price feeds,
        insulating both parties from ETH volatility at the moment
        of commitment.
  \item Records IoT-verifiable logistics milestones on-chain, linking
        each to an IPFS content-identifier (CID) for tamper-evident
        evidence.
  \item Enforces a configurable Delegation-of-Authority (DoA) policy
        that requires $n$-of-$m$ finance approvers before high-value
        funds are released.
  \item Enables on-chain invoice factoring so that suppliers can
        obtain working capital immediately, without waiting for
        delivery confirmation.
\end{itemize}

The remainder of this paper is organised as follows.
Section~\ref{sec:background} surveys related work.
Section~\ref{sec:design} presents the system architecture and protocol
design.
Section~\ref{sec:impl} details the implementation.
Section~\ref{sec:eval} reports gas consumption and test results.
Section~\ref{sec:security} analyses security properties.
Section~\ref{sec:discussion} discusses limitations and future work, and
Section~\ref{sec:conclusion} concludes.

% ============================================================
\section{Background and Related Work}
\label{sec:background}

\subsection{Blockchain in Trade Finance}

The application of distributed ledger technology to trade finance has
been widely studied \cite{guo2019,guo2024,luo2024}.
Pinna and Ruttenberg \cite{pinna2016} identified finality, atomicity,
and programmability as the key properties blockchain brings to
financial settlement.
Consortium platforms such as we.trade \cite{wetrade2022} and Contour
\cite{contour2020} demonstrated feasibility but suffered from high
governance overhead and lack of public accessibility.

\subsection{Decentralised Escrow}

Escrow smart contracts have been proposed for freelance work
\cite{cai2019} and real-estate transactions \cite{vos2017}, yet
few anchor the locked amount to a USD-denominated invoice,
preferring instead to specify ETH directly.
This exposes parties to currency risk: a 20\% ETH price move during
the settlement window can erode half a month's supplier margin.

\subsection{Oracle Integration}

Connecting on-chain contracts to off-chain real-world data (prices,
events) is the oracle problem \cite{zhang2016,caldarelli2020}.
Chainlink decentralised oracle networks provide aggregated,
tamper-resistant price feeds with on-chain freshness metadata
\cite{chainlink2021,chainlinkdocs}.
TracePay leverages the \texttt{AggregatorV3Interface} to obtain a validated
ETH/USD price with strict stale-data detection.

\subsection{Invoice Financing on Blockchain}

Tokenised receivables and on-chain factoring have been explored
by Centrifuge \cite{centrifuge2020}, TradeTrust \cite{tradetrust2020},
and fraud-resilient factoring protocols \cite{guerar2020}.
Our contribution differs in that financing is integrated directly
into the escrow lifecycle rather than operating as a secondary
token market, eliminating the need for a separate ERC-20 token
or bonding curve.

% ============================================================
\section{System Design}
\label{sec:design}

\subsection{Actors and Roles}

TracePay defines four actor types, managed via OpenZeppelin's
Role-Based Access Control (RBAC):

\begin{itemize}
  \item \textbf{Buyer}: Creates and funds orders; may claim a refund
        if delivery is missed.
  \item \textbf{Supplier}: Receives payment upon confirmed delivery;
        may request early financing.
  \item \textbf{Verifier} (\texttt{VERIFIER\_ROLE}): A trusted
        logistics entity (3PL, IoT gateway) that records
        shipment milestones.
  \item \textbf{Approver} (\texttt{APPROVER\_ROLE}): Finance
        controllers who must co-authorise release of high-value orders
        under the DoA policy.
  \item \textbf{Financier} (\texttt{FINANCIER\_ROLE}): Provides
        working capital to the supplier by funding the invoice advance.
\end{itemize}

\subsection{Order Lifecycle State Machine}

The order status follows a strict directed acyclic graph (DAG),
ensuring every transition is validated on-chain:

\begin{align}
  \texttt{Created} &\xrightarrow{\text{fundOrder}} \texttt{Funded} \nonumber\\
  &\xrightarrow{\text{Dispatched}} \texttt{Dispatched} \nonumber\\
  &\xrightarrow{\text{Delivered}} \texttt{Delivered} \nonumber\\
  &\xrightarrow{\text{releasePayment}} \texttt{Settled}
  \label{eq:fsm}
\end{align}

\noindent A parallel refund path is available when the delivery
deadline expires before the \texttt{Delivered} state is reached:

\begin{equation}
  \{\texttt{Funded}, \texttt{Dispatched}\}
  \xrightarrow[\Delta t > \text{deadline}]{\text{claimRefund}}
  \texttt{Refunded}
\end{equation}

\subsection{Oracle-Priced Escrow}

Rather than specifying a fixed ETH amount, buyers specify the invoice
value in USD (8-decimal Chainlink format, e.g., \$1{,}000 is stored as
$1000 \times 10^8$).
At the moment of \texttt{fundOrder}, the contract fetches the live
ETH/USD price $P$ from the Chainlink feed and computes the required
escrow in wei:

\begin{equation}
  w_{\text{escrow}} = \frac{u_{\text{USD}} \times 10^{d_{\text{feed}}} \times 10^{18}}
                           {P \times 10^8}
  \label{eq:usd_to_wei}
\end{equation}

\noindent where $d_{\text{feed}}$ is the feed's decimal count (8 for
all Chainlink USD pairs), $u_{\text{USD}}$ is the invoice USD amount
(also 8-decimal), and $P$ is the raw integer price returned by the
feed.
The price is rejected if: (i) it is $\le 0$, (ii) the round is
incomplete ($\textit{answeredInRound} < \textit{roundId}$), or (iii)
the data age exceeds \texttt{maxPriceAge} (default 2\,hours,
matching the Sepolia ETH/USD heartbeat of 1\,hour).

\subsection{Milestone Verification}

Only addresses holding \texttt{VERIFIER\_ROLE} may call
\texttt{recordMilestone}.
Two milestones are supported: \texttt{Dispatched} (order has left
the warehouse) and \texttt{Delivered} (proof-of-delivery received).
Each call accepts an \texttt{evidenceCid} parameter---a
content-addressed IPFS hash of the supporting document (bill of lading,
IoT telemetry log, e-signature).
This ensures the on-chain log is both tamper-evident and
storage-efficient.

\subsection{Delegation of Authority (DoA)}

For orders whose USD value meets or exceeds the DoA threshold
$\tau_{\text{DoA}}$, at least $k$ distinct approvers must invoke
\texttt{approveRelease} before \texttt{releasePayment} succeeds.
Two invariants are enforced:
\begin{enumerate}
  \item \textbf{Separation of duties}: approvers cannot be the buyer
        or supplier of the order.
  \item \textbf{Uniqueness}: each approver address can grant approval
        at most once per order.
\end{enumerate}

The default configuration sets $\tau_{\text{DoA}} = \$5{,}000$ (8-decimal) and $k = 2$.

\subsection{Payment Waterfall}

Upon successful \texttt{releasePayment}, funds are distributed in
the following priority order (Fig.~\ref{fig:waterfall}):

\begin{enumerate}
  \item Financier receives \texttt{financierDueWei} (advance +
        financing fee), capped at the total escrow.
  \item Supplier receives the remainder.
\end{enumerate}

All credits are written to a \texttt{claimable} pull-payment ledger,
following the \emph{checks-effects-interactions} pattern and
eliminating reentrancy risk on the settlement path.

\begin{figure}[htbp]
  \centering
  \setlength{\unitlength}{1pt}
  \begin{picture}(220,70)
    % Escrow box
    \put(0,10){\framebox(220,55){}}
    \put(5,52){\scriptsize\texttt{escrowWei}}
    % Financier slice
    \put(0,10){\framebox(95,55){}}
    \put(18,34){\footnotesize Financier}
    \put(10,22){\tiny \texttt{financierDueWei}}
    % Supplier slice
    \put(95,10){\framebox(125,55){}}
    \put(132,34){\footnotesize Supplier}
    \put(128,22){\tiny remainder}
  \end{picture}
  \caption{Payment waterfall: financier is paid first (advance + fee),
           supplier receives the balance.}
  \label{fig:waterfall}
\end{figure}

\subsection{Invoice Financing (Factoring)}

A supplier with a funded escrow may call \texttt{requestFinancing}.
A financier holding \texttt{FINANCIER\_ROLE} then calls
\texttt{fundFinancing}, depositing exactly the advance amount
(computed off-chain via \texttt{financingQuote}):

\begin{align}
  w_{\text{advance}} &= \frac{w_{\text{escrow}} \times r_{\text{adv}}}{10{,}000}
  \label{eq:advance}\\
  w_{\text{due}}     &= w_{\text{advance}} \times \left(1 + \frac{r_{\text{fee}}}{10{,}000}\right)
  \label{eq:due}
\end{align}

\noindent where $r_{\text{adv}} = 9{,}000$ (90\%) and
$r_{\text{fee}} = 200$ (2\%) by default, both expressed in
basis-points to avoid floating-point arithmetic.
The supplier immediately receives $w_{\text{advance}}$ in working
capital, while the financier is owed $w_{\text{due}}$ from the escrow
at settlement.

% ============================================================
\section{Implementation}
\label{sec:impl}

\subsection{Technology Stack}

TracePay is built with the following components:
\begin{itemize}
  \item \textbf{Solidity 0.8.24} with \texttt{optimizer: \{runs: 200\}}.
  \item \textbf{OpenZeppelin Contracts v5} --- \texttt{AccessControl},
        \texttt{ReentrancyGuard} \cite{openzeppelin}.
  \item \textbf{Hardhat 2.x} --- compilation, local EVM, testing.
  \item \textbf{Hardhat Gas Reporter} --- per-function gas measurement.
  \item \textbf{MockV3Aggregator} --- Chainlink mock for deterministic
        testing (price and staleness controllable).
  \item \textbf{Ethers.js v6 + MetaMask} --- browser-based frontend
        on Ethereum Sepolia testnet.
\end{itemize}

\subsection{Key Data Structures}

\begin{lstlisting}[caption={TracePay Order struct and state enum.},
                   label={lst:struct}]
enum Status {
  Created, Funded, Dispatched,
  Delivered, Settled, Refunded
}

struct Order {
  address buyer;
  address supplier;
  bytes32 invoiceHash;    // keccak256 hash
  uint256 invoiceUsd;     // 8-decimal USD value
  uint256 escrowWei;      // locked ETH in wei
  uint256 lockedPrice;    // ETH/USD at funding time
  uint64  deadline;       // UNIX delivery deadline
  Status  status;
  uint8   approvals;      // DoA approvals collected
  bool    financeRequested;
  address financier;
  uint256 financierDueWei; // advance + fee owed
}
\end{lstlisting}

\subsection{Oracle Price Validation}

\begin{lstlisting}[caption={Stale-price guard in \texttt{latestPrice()}.},
                   label={lst:oracle}]
function latestPrice()
    public view returns (uint256 price) {
  (uint80 roundId, int256 answer,,
   uint256 updatedAt, uint80 answeredInRound)
      = priceFeed.latestRoundData();
  require(answer > 0, "oracle: price<=0");
  require(updatedAt != 0 &&
          answeredInRound >= roundId,
          "oracle: incomplete round");
  require(block.timestamp - updatedAt
          <= maxPriceAge, "oracle: stale price");
  price = uint256(answer);
}
\end{lstlisting}

\subsection{Settlement and Pull-Payment Pattern}

Following the Ethereum security best-practice of
\emph{checks-effects-interactions} (CEI) \cite{solidity_patterns},
all ETH transfers are deferred to a dedicated \texttt{withdraw()}
function:

\begin{lstlisting}[caption={Pull-payment withdrawal.},
                   label={lst:withdraw}]
function withdraw() external nonReentrant {
  uint256 amt = claimable[msg.sender];
  require(amt > 0, "nothing to withdraw");
  claimable[msg.sender] = 0;      // CEI: effect first
  (bool ok,) = payable(msg.sender)
                 .call{value: amt}("");
  require(ok, "transfer failed");
  emit Withdrawn(msg.sender, amt);
}
\end{lstlisting}

% ============================================================
\section{Evaluation}
\label{sec:eval}

\subsection{Test Suite}

The test suite comprises 16 Mocha/Chai test cases executed against a
Hardhat in-process EVM node.
Table~\ref{tab:tests} summarises the coverage.

\begin{table}[htbp]
  \caption{TracePay Test Coverage Summary}
  \label{tab:tests}
  \centering
  \small
  \renewcommand{\arraystretch}{1.15}
  \begin{tabular}{cp{5.3cm}c}
    \toprule
    \textbf{TC} & \textbf{Scenario} & \textbf{Result} \\
    \midrule
    TC-01 & Happy path: create $\to$ fund $\to$ deliver $\to$ settle & \checkmark \\
    TC-02 & Oracle conversion correctness at varying ETH prices       & \checkmark \\
    TC-03 & DoA: high-value order requires 2 distinct approvals        & \checkmark \\
    TC-04 & Invoice financing waterfall (advance + fee)                & \checkmark \\
    TC-05 & Refund after deadline when undelivered                     & \checkmark \\
    TC-06 & Overpayment credited back to buyer via pull-payment        & \checkmark \\
    TC-07 & Duplicate invoice hash rejected                            & \checkmark \\
    TC-08 & Stale oracle price blocks funding                          & \checkmark \\
    TC-09 & Negative oracle price blocks quoting                       & \checkmark \\
    TC-10 & Underfunding ETH reverts                                   & \checkmark \\
    TC-11 & Only verifier records milestones in order                  & \checkmark \\
    TC-12 & Separation of duties: buyer cannot approve release         & \checkmark \\
    TC-13 & Same approver cannot approve twice                         & \checkmark \\
    TC-14 & Premature release and double-release both revert           & \checkmark \\
    TC-15 & Non-financier and wrong-amount financing revert            & \checkmark \\
    TC-16 & Empty-balance withdraw and invalid refund revert           & \checkmark \\
    \midrule
    \multicolumn{2}{l}{\textbf{Total passing test cases}} & \textbf{16 / 16} \\
    \bottomrule
  \end{tabular}
\end{table}

All 16 tests pass deterministically on a local EVM.

\subsection{Gas Consumption}

Table~\ref{tab:gas} reports gas measurements obtained with Hardhat
Gas Reporter (Solc 0.8.24, optimiser enabled, 200 runs,
block limit 30,000,000).

\begin{table}[htbp]
  \caption{Gas Consumption by Function}
  \label{tab:gas}
  \centering
  \small
  \renewcommand{\arraystretch}{1.12}
  \begin{tabular}{lrrr}
    \toprule
    \textbf{Function} & \textbf{Min} & \textbf{Max} & \textbf{Avg} \\
    \midrule
    \texttt{createOrder}    & \multicolumn{2}{c}{181,063} & 181,063 \\
    \texttt{fundOrder}      & 94,510  & 116,781 &  96,535 \\
    \texttt{fundFinancing}  & \multicolumn{2}{c}{85,353}  &  85,353 \\
    \texttt{releasePayment} & 61,888  &  69,569 &  64,333 \\
    \texttt{approveRelease} & \multicolumn{2}{c}{58,259}  &  58,259 \\
    \texttt{claimRefund}    & \multicolumn{2}{c}{57,556}  &  57,556 \\
    \texttt{grantRole}      & 51,481  &  51,493 &  51,490 \\
    \texttt{recordMilestone}& 32,984  &  32,995 &  32,990 \\
    \texttt{withdraw}       & \multicolumn{2}{c}{32,410}  &  32,410 \\
    \texttt{requestFinancing}& \multicolumn{2}{c}{30,628} &  30,628 \\
    \midrule
    \textit{Deployment}     & 2,227,080 & 2,227,092 & 2,227,091 \\
    \textit{(7.4\% of block limit)} & & & \\
    \bottomrule
  \end{tabular}
\end{table}

The most expensive operation is \texttt{createOrder} (181,063 gas)
due to storage initialisation.
All subsequent lifecycle operations require less than 100,000 gas units.

\subsection{Deployment Footprint}

The contract deploys in 2,227,091 gas units on average, occupying 7.4\%
of the 30,000,000-gas Ethereum block limit, ensuring rapid inclusion
even during periods of moderate network congestion.

% ============================================================
\section{Security Analysis}
\label{sec:security}

\subsection{Reentrancy}

All state-mutating external calls follow the CEI pattern, and the
\texttt{nonReentrant} modifier from OpenZeppelin's
\texttt{ReentrancyGuard} is applied to every function that
moves ETH (\texttt{fundOrder}, \texttt{fundFinancing},
\texttt{releasePayment}, \texttt{claimRefund}, \texttt{withdraw}).
The pull-payment model further decouples settlement logic from
ETH transfer.

\subsection{Access Control}

Three custom roles plus the default admin role are enforced by
OpenZeppelin \texttt{AccessControl}.
Table~\ref{tab:rbac} maps each privileged function to its required role.

\begin{table}[htbp]
  \caption{Role-to-Function Access Map}
  \label{tab:rbac}
  \centering
  \small
  \renewcommand{\arraystretch}{1.12}
  \begin{tabular}{lp{4.2cm}}
    \toprule
    \textbf{Role} & \textbf{Functions} \\
    \midrule
    \texttt{VERIFIER\_ROLE}  & \texttt{recordMilestone} \\
    \texttt{APPROVER\_ROLE}  & \texttt{approveRelease} \\
    \texttt{FINANCIER\_ROLE} & \texttt{fundFinancing} \\
    \texttt{DEFAULT\_ADMIN}  & \texttt{setPolicy}, \texttt{setFinancingTerms}, \newline \texttt{setMaxPriceAge} \\
    \bottomrule
  \end{tabular}
\end{table}

\subsection{Oracle Manipulation}

The Chainlink ETH/USD feed is decentralised and aggregated
from multiple node operators.
TracePay adds a defence-in-depth layer: the price is locked at the
moment of \texttt{fundOrder} and stored in \texttt{lockedPrice} for
auditability.
A post-funding price change cannot affect an already-committed escrow
amount.

\subsection{Front-Running and Miner Extractable Value (MEV)}

The \texttt{fundOrder} function is restricted to the order's
buyer.
An adversarial observer cannot substitute their own address or
divert the locked funds.
The oracle freshness window (2\,hours) prevents exploitation of
outdated round prices.

\subsection{Integer Arithmetic}

Solidity 0.8.x includes native checked arithmetic preventing overflows.
All basis-point calculations use 256-bit unsigned integers without
precision loss.

\subsection{Duplicate Invoice Prevention}

The \texttt{invoiceRegistered} mapping records every accepted
\texttt{invoiceHash} (keccak256 hash of the off-chain invoice).
Replay attempts immediately revert with \texttt{"duplicate invoice"}.

% ============================================================
\section{Discussion and Future Work}
\label{sec:discussion}

\subsection{Limitations}

\textbf{Oracle dependency.}
Although Chainlink mitigates single points of failure, reliance on
external feeds introduces residual risk during extreme network partitions.
A multi-oracle router could provide secondary fallback sources.

\textbf{Verifier centralization.}
Milestone recording relies on authorized verifiers holding
\texttt{VERIFIER\_ROLE}.
Integrating cryptographically signed IoT telemetry or zero-knowledge
proofs can reduce human verifier dependency.

\textbf{Volatility during escrow.}
Because escrow is locked in ETH, intra-shipment currency movement
impacts the fiat-equivalent value received at settlement.
Integrating ERC-20 stablecoins (e.g., USDC) provides price stability.

\subsection{Future Work}

Future enhancements include:
\begin{itemize}
  \item Multicurrency and stablecoin settlement (USDC / ERC-20).
  \item NFT-tokenized receivables (ERC-721) for secondary market trading.
  \item Zero-knowledge proofs for confidential milestone verification.
  \item Layer-2 rollups (Arbitrum, Base) deployment for 10--100$\times$ gas reduction.
\end{itemize}

% ============================================================
\section{Conclusion}
\label{sec:conclusion}

TracePay demonstrates a trust-minimized trade finance platform
encapsulated in a single Solidity smart contract.
By combining oracle-priced escrow, multi-signature approval, milestone
verification, and on-chain invoice factoring, TracePay eliminates
intermediary delays and high fees.
With all 16 security and functional tests passing and deployment gas
within 7.4\% of the block limit, TracePay offers a practical,
auditable foundation for decentralized trade finance.

% ============================================================
\balance
\bibliographystyle{IEEEtran}
\bibliography{references}

\end{document}
